\chapter{Conclusions and Future Directions}\label{ch:conclusion}

This thesis presented Motivo, a domain-specific language, native compiler, and web-based studio for offline symbolic
music composition.
The project began with a practical observation: repeated musical structures are often easier to describe as rules than
as copied note events.
Motivo turns that observation into a compiler pipeline.
The source program names musical values and procedures, the semantic analyzer validates the structure, the lowerer
expands it into timestamped events, and the MIDI writer serializes the result into a Standard MIDI File.

The work also demonstrates that the compiler can be embedded into a larger software product.
Motivo Studio connects the C++ compiler to a browser editor through an Express API, renders generated MIDI in a piano
roll, supports browser playback, and is packaged through Docker and nginx.
The project is therefore not only a small programming language experiment.
It is also a full-stack engineering exercise with component boundaries, automated tests, CI workflows, deployment
configuration, and a user-facing interface.

\section{Contributions}\label{sec:contributions}

The first contribution is the Motivo language model.
It provides a compact set of abstractions for notes, rests, chords, sequences, patterns, tracks, voices, and control
flow.
These abstractions are deliberately close to symbolic musical structure and deliberately far from raw MIDI byte
construction.
The language is small, but it is sufficient to express repeated arrangements and parallel lines in the current
implementation.

The second contribution is the compiler architecture.
The implementation uses a conventional staged pipeline: scanner, parser, semantic analyzer, lowerer, and backend.
The most project-specific part is the cursor-based lowering model.
That model lets sequential statements advance musical time, explicit offsets place events without changing the main
cursor, voices run with isolated cursors, and pattern calls return reconstructed musical values.
This is the core mechanism that connects source-level structure to MIDI note events.

The third contribution is Motivo Studio.
The application shows how the compiler can be made accessible through a web interface.
The editor, diagnostics loop, MIDI parsing, piano-roll visualization, playback controls, server wrapper, Docker stack,
and CI workflows demonstrate the software engineering side of the project.
For a bachelor's thesis, this part is important because it shows that the work is not only theoretical.
It is an implemented system with several interacting modules.

The same contribution can also be viewed through the structural abstractions implemented by the compiler. The
language provides constructs such as \texttt{pattern}, \texttt{track}, and \texttt{voice}, which allow
polyphonic arrangements to be described without manually calculating absolute timestamps or delta-times. The
deterministic lowering algorithm connects those abstractions to a static intermediate representation, while
the semantic analyzer checks type and structural constraints before MIDI generation.

The fourth contribution is the empirical evaluation.
The local benchmark results show that Motivo can generate large symbolic event streams from compact source programs and
compile them quickly in the tested Release configuration.
The results are intentionally bounded by the current implementation and hardware.
They support the practical claim that Motivo is suitable for interactive local experimentation at the current scale.

\section{Limitations}\label{sec:limitations}

Motivo is not a complete composition environment.
It does not provide audio synthesis inside the compiler.
It does not model performer expression beyond the simple event fields currently present in the intermediate
representation.
It does not support automation curves, pitch bend, continuous controller data, microtonality, or symbolic harmonic
analysis.
It also does not currently support key-signature declarations.
The language emits MIDI that can be played or imported elsewhere, but further sound design and mixing remain outside
the compiler.

The compiler also has an explicit 20,000-note-event safety limit.
This limit is appropriate for preventing accidental runaway programs in the current implementation, but it constrains
large generated scores.
The benchmark results should therefore be understood within this cap.
Future work could replace the fixed cap with configurable resource limits and better reporting, but the current thesis
documents the implemented behavior as it exists.

The evaluation is technical rather than user-centered.
It measures compilation time, source-to-event expansion, MIDI size, and test evidence.
It does not measure whether composers prefer Motivo, whether beginners understand the syntax more easily than other
systems, or whether pieces written in Motivo are musically better.
Those would require a different study design and participant data.
This thesis avoids those claims.

The current pattern system is also intentionally static. Pattern parameters accept concrete evaluated types,
but the language does not yet support higher-order pattern manipulation, richer harmonic types, or a standard
library of reusable musical transformations. The system is therefore expressive enough for the demonstrated
examples, but it remains a compact language rather than a complete compositional platform.

\section{Future Work}\label{sec:future-work}

The most natural language extension is richer musical context.
Key signatures are a good example because they would let Motivo understand scale degrees, diatonic transposition, and
harmonic functions instead of treating all notes as independent pitch literals.
This should be introduced carefully.
It would require syntax, semantic rules, interactions with accidentals, and clear lowering behavior.
Because that support is not part of the current implementation, it belongs in future work rather than in the language
description.

Another useful extension is continuous musical data.
MIDI control changes, pitch bend, and automation curves would make Motivo more expressive for modern production
workflows.
Supporting them would require changes to the AST, type system, intermediate representation, MIDI backend, piano-roll
visualizer, and playback scheduler.
The existing separation between IR and backend makes such an extension plausible, but it is not a small addition.

The compiler could also expose more evaluation and debugging data.
For example, Motivo Studio could show the number of generated events, the longest track duration, compile time, MIDI
file size, and event-density warnings after each compile.
This would turn the current thesis measurements into user-facing feedback and would help users understand when a source
program is becoming too large.

Finally, the application could support project persistence and export workflows.
Saving Motivo source files, exporting MIDI with metadata, comparing generated versions, and integrating with cloud
storage would move Motivo Studio closer to a practical composition tool.
Those features would not change the compiler's core idea, but they would improve the surrounding product.

A standard library is another natural direction. Reusable scales, chord progressions, arpeggiators, Euclidean
rhythms, and common drum patterns would make the language more practical without changing its core compiler
model. Such a library should be designed after the core type system is stable, because library functions
would depend on how Motivo represents harmonic context and pattern parameters.

Alternative backends also remain possible future work. The current backend targets Standard MIDI Files, but
the same intermediate representation could be used as the basis for JSON event export, MusicXML-oriented
notation export, or integration with synthesis systems. Those targets would require separate mapping rules
and should not be treated as implemented behavior in the present thesis.

\section{Closing Remarks}\label{sec:closing}

Motivo demonstrates that compiler design can be applied to a creative symbolic domain in a concrete and measurable way.
The language gives repeated musical ideas a compact textual form, and the implementation turns that text into MIDI
events through a staged compiler.
The Studio application then places the compiler inside a broader engineering system with editing, feedback,
visualization, playback, deployment, and tests.

The result is not a replacement for existing music tools.
It is a focused bridge between algorithmic thinking and standard MIDI-based production workflows.
That focus is the main strength of the project: it keeps the language small enough to implement and evaluate, while
still showing how software engineering can support structured musical composition.
